\subsection{Expanded case interpretations}
\subsubsection{Specification inflation}\label{sec:m1}

\emph{Interpretation: specification growth outpaced the behaviour validated
through product use.} By 22~October the BRD had reached version~0.19 and
about 37,900 words, with requirement identifiers up to FR-052 and AC-103; the
PRD had reached version~0.18 and about 33,900 words, with identifiers up to
FR-061 and AC-110. The backlog inherited that scale and grew to 192 issues.
Cross-cutting concerns were appended to existing issues rather than sequenced
as their own work: a responsible-AI section on 15~October and an
AI-consumption section on 22~October. Issue \#126, on pilot release gates,
ended with seventeen acceptance criteria in three sections. Even the request
for the smallest runnable preview arrived with thirteen.

This expansion began before the autonomous backlog loop, including after participant
prototype and MVP feedback. The requirements agents sought fully specified
BRDs and PRDs, named approvers and resolved ownership before building. The author
requested explicit demonstrator assumptions and imperfect increments that
could produce feedback. Smaller, iterative requirements documents are a
plausible alternative, not an intervention tested in this case.

The pattern resembles B\"ockeler's small bug that became sixteen acceptance
criteria~\cite{boeckeler2025sdd}, at the scale of a programme, and the risks
of late integration described by Royce and Gall~\cite{royce1970,gall1975}.
Brooks locates essential difficulty in the conceptual construct rather than
its representation~\cite{brooks1987silver}; repeating a requirement in more
places need not resolve its ambiguity. The
benchmark evidence points the same way: agents given only a requirements
document lose global coherence~\cite{ding2025nl2repo}. Beck's metaphor fits:
features consume future options~\cite{beck2026prodacity,beck2025devslow}.
Extending that argument to acceptance criteria is the present author's
interpretation, not Beck's wording. The early-warning signal is
the ratio of specification growth to behaviour a human has validated in the
same period. No replacement journey was reported complete during the loop,
but a systematic count of validated behaviours was not collected; this ratio
cannot be estimated from the document counts alone.

\subsubsection{Proxy capture}\label{sec:m2}

\emph{In this case, artefacts made into gates appeared to displace outcomes.}
Three patterns motivate this interpretation.

Merged pull requests are a weak proxy for delivered value. Murphy-Hill et~al.
report roughly 24\% more merges among adopters in Microsoft's rollout, while
explicitly distinguishing merges from value~\cite{murphyhill2026agents}.
Peralta et~al. studied 9,799 human-reviewed agent pull requests; in their
717-PR manually inspected sample, only 35.7\% of rejections reflected clear agent
failures and 15.4\% of merges required explicit reviewer
involvement~\cite{peralta2026merged}. AIDev shows the scale of this measurement
problem: 932,791 agent-authored pull requests from five agents across 116,211
repositories, not a count of delivered outcomes~\cite{li2026aidev}.

\emph{Process.} When the first waves of sessions skipped the research and
planning phases, the author made the presence of their artefacts a merge condition.
From then on every pull request carried research, plan, change and review
records; the artefacts were guaranteed, the thinking they were meant to
evidence was not. Pull request \#445 changed two lines of behaviour and added
a 78-line evidence record that reports its own artefact completion as 100\%.
HVE Core treats these files as working memory. Substantial tracking output
was already committed before the added merge condition; that condition
formalised an existing practice rather than explaining its onset.

In the first-parent integration view, tracking-artefact additions were
10.27 times product-source additions during the loop, compared with 29.28
before it (\nolinkurl{anc/repo_history_metrics.out.txt}). Artefact volume remained
high during the loop, while its ratio to product-source additions fell.
Churn alone does not establish code decay.

\emph{Evaluation.} The rubric scorer, a language model reading the repository,
rated implementation completeness at five out of five, citing a README
and a check script, at a
time when the author could not sign in to the replacement. The score and blocked
sign-in occurred during the same period. It remains necessary to check whether the
score covered that release and journey (Section~\ref{sec:audittrail}).
The judge read repository text rather
than executing the intended journey. Neither its repeatability nor its
agreement with observed product behaviour was measured; its score cannot
substitute for that behavioural check.

\emph{Evidence boundaries.} The author had instructed the agents never to
fabricate approvals or overstate evidence. The record contains explicit
statements of unmet acceptance, and lexical evidence-boundary density more than tripled.
These observations do not establish exhaustive compliance or measure honesty.
Issue \#349 records that a browser redirect ``is not successful sign-in''.
The agents were right, and still nobody could sign in. Honest boundary
statements became the output.

Goodhart and Strathern predict the first two
levels~\cite{goodhart1975,strathern1997improving}. Beck joked that an agent
seemed ready to delete tests and pretend it had finished
~\cite{orosz2025beck}; ImpossibleBench measures test exploitation more
generally, not a counted test-deletion category~\cite{zhong2026impossiblebench}.
The METR study adds the human side: the
feeling of progress is an unreliable instrument~\cite{becker2025metr}. The
early-warning signals are the ratio of evidence artefacts to behaviour
changes, the share of merged pull requests that change what a user can do,
boundary-phrase density, and any divergence between automated scores and a live demo.

\subsubsection{Learning evaporation}\label{sec:m3}

\emph{Some recorded lessons appeared to fall off the loop's read path, while
important design decisions survived.} Maintaining the structured log and
memory became difficult as the record grew. Important design decisions
survived, even as some recorded lessons appeared to be lost. Before the loop, participant
use of the MVP produced concrete lessons: the simplicity requirements, a
deliberately public development profile for throwaway stacks, and the
callback, role and staff-binding steps that made sign-in work. All were
recorded, in the PRD, in a simplicity review and in infrastructure runbooks.
The stated brief format carried issue text, acceptance criteria and
repository conventions; actual session access to the earlier lessons is
unverified. The November search finds retained Entra, registration,
staff-binding and runbook references in committed research/plans,
including 16 October records. This weakens a blanket disappearance
claim without establishing that agents accessed or used the material. The log cites the
simplicity requirements for the last time on 13~October; the development
profile appears only on 8~October. The migrated deployment lacked the
registration and staff bindings established on 12~October, and sign-in was
blocked.

The read path grew. The knowledge-base index, which every agent is
told to read first, accumulated 1,069 words of issue-specific status notes
ahead of its first instruction. The index, the log, the BRD and the PRD, the
documents that define the project's state and intent, total about 139,000
words. No provider-matched tokenizer or actual session-input record was
available; prose token-per-word heuristics are unsuitable for these
identifiers, paths and tables.
Long contexts can reduce reliable retrieval even when the material fits in
the context window~\cite{liu2024lost,hong2025contextrot}; the case does not
measure any session's effective attention capacity.

Instruction debt and learning evaporation can therefore coexist. A correction
may remain in an expanding agent definition while the evidence explaining
when it applies falls off the next session's read path. The result can be
a larger instruction set with weaker access to the relevant lesson.
However, unavailable guidance, accessed-but-misapplied guidance and guidance
insufficient for the task are different explanations
\cite[Secs.~5.2--5.5]{jiang2026skills}. Without actual session inputs and
access/use traces, this case cannot distinguish them.

Naur explains why recording was not enough. A program's theory lives in the
people who build it, and a second team with full documentation can still
destroy a design's simplicity~\cite{naur1985theory}. In an agent loop, every
session may face an analogous reconstruction problem. The comparison with
Naur is interpretive. The case provides no measure of lost understanding.
Yan's advice to share context appears within an argument against
multi-agent designs~\cite{yan2025multiagents}, not evidence that adding
agents solves retrieval. The multi-agent failure taxonomy
~\cite{cemri2025multiagent} identifies several coordination problems; it
does not establish context loss as a prevalent cause in this case. HVE
Core's choice to carry context in files on disk is necessary but not
sufficient: the files must be on the read path, short enough to read and,
ideally, executable. The early-warning signals are lessons that exist only as
prose, a read path that grows faster than the code, and problems that have to
be solved twice.

Pavlenko et~al.'s World Model offers a different design: preserve raw evidence,
derive rebuildable views and retrieve compact, task-specific context with
provenance and access controls. Updates are reviewed, and measured outcomes
feed later decisions~\cite{pavlenko2026factory}. This distinguishes storing
lessons from making them usable in the next decision; a small read path need
not mean discarding history. This architecture offers another way to retrieve
accumulated lessons. Its effect on learning evaporation in this case remains
untested.

\subsubsection{Process-gate inflation}\label{sec:m4}

\emph{Interpretation: repeated gates appeared disproportionately burdensome;
comparative effort was not measured.} Each new gate required execution and
review. Within ten days the coordinator required, for every
issue, the full RPI loop, a complete local gate before every push and after
every correction or rebase, independent review, up to two correction cycles,
board updates and evidence-based ticking of acceptance criteria. Hosted CI
minutes had to be rationed, the full gate had to be interrupted by explicit
exception for a two-line fix, and paid cloud validation stayed blocked even
after the cost ceiling was raised.

The author perceived attention as more constrained than compute, but their
relative scarcity was not measured. Reviewing eighteen merged pull requests
a day in depth was judged impractical, and merging was delegated to the
coordinator. The author still answered approvals and technical blockers,
but stopped getting feedback from using a replacement product. RPI critiques
and reviews addressed the current issue rather than the cumulative system;
there was no regular phase for reflective consolidation or generalisation
across completed issues. This is Bainbridge's irony in
miniature~\cite{bainbridge1983ironies}: the automation took the part of the
work that supplied direct product-use feedback and left the participant
monitoring a stream of plausible, well-documented changes. It
matches the shift from producing to evaluating reported for generative
AI~\cite{simkute2024ironies}, the delegation demands it
creates~\cite{tankelevitch2024metacognitive}. Higher confidence in AI was
associated with reduced self-reported critical scrutiny in the cited
study~\cite{lee2025critical}, not established as its cause. Murphy's
diagnosis is a plausible interpretation: throughput increased while the path
to a usable outcome remained blocked~\cite{murphy2026speed,goldratt1984goal}.
The bottleneck itself was not instrumented. Martin's argument for executable
feedback from tests and code-analysis tools~\cite{martin2026productengineers}
motivates a proposed gate that checks product behaviour rather than inferring
it from paperwork; the interview does not establish that such a gate would
have prevented this failure. The early-warning signals are gate time per change,
human review minutes per merged pull request, work in progress, and cost per
merged pull request compared with cost per accepted outcome.

\subsubsection{Feedback starvation}\label{sec:m5}

\emph{Candidate mechanism: insufficient independent, outcome-level feedback
even though artefact-level validation is abundant.} This differs from merely having
few tests, a long feedback delay or missing context: many checks can pass
while none disconfirms the claim that the intended user journey works.
Outcome-level independence here means evidence obtained by exercising the outcome rather
than rereading its completion record; it need not imply a different model.
It does not establish independent human judgment.
The loop prioritised completing the specified plan rather than shortening the
path to such feedback. Feedback starvation is a central interpretive proposition,
although the case cannot separate its contribution from the other mechanisms.
Nothing in the coordinator kept a usable system alive:
no invariant required the participant's sign-in and core workflow to remain
usable after every merge. Completion was assessed through artefacts, criteria
and gates. In the MVP week the feedback loop included intermittent participant use, with some changes made the same
day. From 15~October the first priority was restructuring code, the
working MVP was frozen, and its replacement had not become usable by the time
of the case cut-off. When a usable increment was
requested, the request was absorbed by the same
machinery: a label, new issues, thirteen acceptance criteria and the full set
of gates.

Several older practices in Section~\ref{sup:sec:theory} make feedback and working
behaviour explicit: short releases and test-first development~\cite{beck2004xp,beck2002tdd},
incremental delivery in Crystal Clear~\cite{cockburn2004crystal},
tracer bullets~\cite{hunt1999pragmatic}, and working software
over comprehensive documentation~\cite{agile2001manifesto}. Beck's warning
that fully agentic work without feedback goes off the
rails~\cite{beck2026prodacity}, DORA's emphasis on small batches and user
focus~\cite{dora2025} and Murphy's list of real bottlenecks, which includes
missing feedback~\cite{murphy2026speed}, motivate this interpretation but do
not independently diagnose the case. The
early-warning signals are the time from the start of a cycle to the first
human use of a running increment, and the age of the last increment a human
has verified as usable. The case records no completed replacement journey
during the ten-day loop; the exact last-use timestamp remains unverified.

\subsection{Case rival-explanation register}\label{app:rivals}

The main interpretation is in Section~\ref{sec:threats}. The register preserves
recorded rivals separately from untested alternatives. A limit is not grounds
for ruling an explanation out. Score/deployment correspondence remains
unverified (Section~\ref{sec:audittrail}).

\begin{table}[!htbp]
\centering\footnotesize
\caption{Rival explanations with direct case records.}\label{tab:rivals}
\begin{tabularx}{\textwidth}{@{}L{2.7cm}YYY@{}}
\toprule
Explanation & Case support & Counter-evidence or limits & Treatment \\
\midrule
Migration decision & ADR 0006 (issue \#196) froze feature work in the MVP and
prioritised a new layout. & Health and restart checks passed, not usability. &
Migration and coordination changed together; autonomy alone is not identified. \\
Capstone rubric & Research and documentation are rewarded; full completeness
marks coexist with blocked sign-in. & A working demonstrator is required;
release/scope linkage is pending. & Possible incentive, not a measured effect. \\
Owner overrides and priorities & Committed RPI records, full gates and delegated
merging; PR \#445 required an exception. & The MVP used agents successfully;
HVE Core defaults retain human stops. & Owner choices are part of the case. \\
Model and tool changes & Versions changed during the study. & No matched rerun
under fixed versions. & Actual session configuration remains partly unknown. \\
CI and cloud budgets & Issue \#413 moved validation locally; costs prompted
a higher ceiling and reduced resources. & Merges continued without established
deployed usability. & More budget is not proven to resolve delivery. \\
Identity and sign-in setup & Missing registration/staff bindings; issue \#349
distinguishes redirect from sign-in. & Those steps worked on 12~October. &
Concrete configuration obstacle, not proof of code decay or lost knowledge. \\
\bottomrule
\end{tabularx}
\end{table}

\FloatBarrier
\begin{table}[!htbp]
\centering\footnotesize
\caption{Additional plausible rivals and unperformed discriminating checks.}\label{tab:rivalsadditional}
\begin{tabularx}{\textwidth}{@{}L{2.7cm}YYY@{}}
\toprule
Explanation & Reason considered & Counter-evidence or limits & Evidence still needed \\
\midrule
External services and deployment drift & Model/identity configuration differs
from the MVP. & No provider outage is established; health checks were narrower. &
Versioned deployment/configuration and dependency-health records for the journey. \\
Scope expansion & RAI and consumption criteria added to issues; further
resident-preview obligations. & Expansion predates autonomy; some controls
are justified. & Original/added scope, sequencing and time-to-outcome comparison. \\
Inadequate acceptance tests & Checks passed without completed real-user sign-in;
scorer read artefacts. & Bounded, honest tests need not be poor tests. &
Capability-to-task-oracle mapping and execution evidence. \\
Rate and context limits & Parallel sessions and a large read path. & No runtime
rate-limit, truncation or attention dataset analysed. & Actual inputs and
blocked/retried request events. \\
Specification ambiguity & Selective review of long, cross-referenced requirements.
& Important domain/consent decisions survived; ambiguity not independently coded.
& Independent interpretation and conflicting acceptance examples. \\
Ordinary prototype debt & Rapid MVP followed by restructuring and rebuilding.
& MVP worked intermittently; no comparable human-led migration. & Failure
analysis distinguishing inherited debt from introduced faults. \\
\bottomrule
\end{tabularx}
\end{table}
\FloatBarrier

\subsection{Candidate-mechanism evidence audits}\label{app:mechanismaudit}

These are operational boundaries for this paper, not validated necessary and
sufficient conditions for all agentic development. The indicators refer to
Table~\ref{tab:audittrail} and the case chronology. Experimental implications
state what evidence would discriminate an explanation; the tagged experiment
does not manipulate the five candidates separately.

\subsubsection{Specification inflation}

{\small
\noindent
\begin{tabularx}{\textwidth}{@{}L{2.6cm}Y@{}}
\toprule
Element & Audit \\
\midrule
Definition & Requirements and acceptance obligations expand faster than
independently validated behaviour, without evidence that the added scope is
needed for the next usable increment. Length alone is insufficient. \\
Case indicators & O3: BRD/PRD growth, 192 issues and cross-cutting additions;
the smallest resident preview had thirteen criteria. \\
Counter-evidence & Growth began before autonomy and sometimes followed the
owner's feedback; regulatory and consent controls can justify detail. \\
Rival explanations & Genuine scope expansion, specification ambiguity,
capstone incentives and an ambitious migration rather than agent-driven
inflation. \\
Experimental implication & Compare smaller requirements horizons at matched
scope and resources. The current fixed-suite experiment cannot identify this
mechanism, even if consolidation helps. \\
\bottomrule
\end{tabularx}
\par
}

\subsubsection{Proxy capture}

{\small
\noindent
\begin{tabularx}{\textwidth}{@{}L{2.6cm}Y@{}}
\toprule
Element & Audit \\
\midrule
Definition & A completion proxy becomes a decision criterion and remains
positive when an independent test of the intended outcome is negative. Proxy
use alone does not establish capture. \\
Case indicators & O2--O3: a 5/5 score and blocked sign-in coexist in the
chronology; release/scope linkage is pending. Mandatory artefact presence and
the evidence record for PR \#445 are separate indicators. \\
Counter-evidence & Agents explicitly reported unmet acceptance; neither
fabricated approval nor deliberate deception is established. \\
Rival explanations & An inadequate scorer or acceptance oracle, different
evaluation targets, and a deployment failure invisible to repository review. \\
Experimental implication & Disagreement between visible completion and hidden
behavioural tests would expose proxy limits. It would not by itself establish
intentional test gaming or explain why a gate was chosen. \\
\bottomrule
\end{tabularx}
\par
}

\subsubsection{Learning evaporation}

{\small
\noindent
\begin{tabularx}{\textwidth}{@{}L{2.6cm}Y@{}}
\toprule
Element & Audit \\
\midrule
Definition & A recorded lesson relevant to a later decision is absent from,
unusable on, or not applied through that decision's effective context path,
and the associated problem recurs. A missing citation alone is insufficient. \\
Case indicators & O4: simplicity references cease; working sign-in setup is
not preserved; state and intent documents reach about 139,000 words. \\
Counter-evidence & Important domain decisions survived; the lesson could have
been known but deferred. Actual session attention was not measured. \\
Rival explanations & Deliberate migration sequencing, deployment-specific
configuration, context limits, or conflicting requirements rather than lost
knowledge. \\
Experimental implication & Trace actual inputs and repeated faults. B--C
cannot isolate memory quality because DESIGN and REVIEW notes and acting
permissions differ. A matched-context study would be needed. \\
\bottomrule
\end{tabularx}
\par
}

\subsubsection{Process-gate inflation}

{\small
\noindent
\begin{tabularx}{\textwidth}{@{}L{2.6cm}Y@{}}
\toprule
Element & Audit \\
\midrule
Definition & Additional procedural obligations increase per-change burden
without demonstrated proportional improvement in the targeted outcome. More
gates alone are not evidence that their net benefit is negative. \\
Case indicators & O5: full RPI and validation on each correction/rebase;
hosted-minute rationing; explicit exception for PR \#445. \\
Counter-evidence & Checks caught bounded defects and enforced evidence
integrity. Gate duration and human review effort were not measured per PR. \\
Rival explanations & Legitimate safety and privacy requirements, budget
constraints, new deployment complexity and owner-selected risk tolerance. \\
Experimental implication & Measure gate time and accepted outcomes under
matched risk. The pilot's resource description concerns the package, not the effect of
removing this repository's required gates. \\
\bottomrule
\end{tabularx}
\par
}

\subsubsection{Feedback starvation}

{\small
\noindent
\begin{tabularx}{\textwidth}{@{}L{2.6cm}Y@{}}
\toprule
Element & Audit \\
\midrule
Definition & Independent outcome-level feedback is insufficient to correct
delivery decisions despite abundant artefact-level validation. The defining
gap is between the evaluated proxy and intended outcome, not test count alone. \\
Case indicators & O1--O2 and O5: working MVP frozen; replacement unusable;
passing checks and continued approvals did not restore a completed user task. \\
Counter-evidence & Human oversight continued; agents surfaced blockers;
engineering checks provided real, though narrower, feedback. \\
Rival explanations & Identity and service configuration, migration priority,
incomplete acceptance coverage, prototype debt or insufficient resources. \\
Experimental implication & Hidden retention tests and CR1 supply behavioural
feedback to the evaluator. The experiment does not test live deployment
recovery, product selection or a human-use checkpoint cadence. \\
\bottomrule
\end{tabularx}
\par
}

\subsection{Measures and lexicon}\label{app:measures}

\paragraph{Log entries.} The knowledge-base log
(\texttt{docs/reference/log.md}) was split into entries at each second-level
heading that begins with a date. An entry belongs to the date in its heading.
Word counts use whitespace tokenisation of the entry, heading included.

\paragraph{Evidence-boundary phrases.} Each entry was searched, case-insensitively and
on word boundaries, for these patterns: \emph{unmet}; \emph{not accepted};
\emph{not inferred}; \emph{is inferred}; \emph{remain(s) open};
\emph{remain(s) unproved}; \emph{remain(s) unverified}; \emph{not claimed};
\emph{not established}; and \emph{no} followed, within forty letters, spaces,
commas, slashes or hyphens, by \emph{is} or \emph{was} and then
\emph{claimed} or \emph{inferred}. Density is matches per 1,000 words.

\paragraph{Reference tracing.} An entry cites the simplicity requirements if
it contains an identifier from SA-001 to SA-010 or the word
\emph{simplicity}. It cites the development profile if it contains
\emph{deploymentProfile}, \emph{development profile} or
\emph{allowPublicNetworkAccess}.

\paragraph{Repository counts.} Issue and pull-request counts come from the
GitHub search API on 24 October 2026 (for example,
\texttt{is:pr is:merged merged:>=2026-10-15}). Document sizes are whitespace
word counts of the files at commit \texttt{29c6b51}.

\paragraph{Proposed measures.} \emph{Time to first human use}: elapsed time
from the start of a delivery cycle to the first time a person uses the running
increment for its purpose. \emph{Age of the last usable increment}: time since
a person last verified that the deployed system could be used for its main
journey. \emph{Behaviour-changing merge ratio}: the share of merged pull
requests that change what a user can observe in the running increment.
\emph{Cost per accepted outcome}: total model, tool and cloud cost divided by
the number of outcomes a human accepted in the same period.

\subsection{Consolidation step prompt}\label{app:prompt}

Verbatim frozen v4 prompt from execution source
\nolinkurl{e8707ecc8965d3f81f5142a2603042bfca1aab15}, not the amended
study's approved prompt. The archived source is byte-identical, SHA-256
\nolinkurl{75f646edbe04e1e3403b092b13edd6b66615a3362e67ad49744d229cdcb94f29}.
Display wrapping does not alter its text. Braces mark runner parameters.

\VerbatimInput[breaklines=true,breakanywhere=true,fontsize=\footnotesize]
{anc/consolidation-frozen-v4.prompt.txt}
